\section{Related Work}

Previous work on differential testing is considerable, but previous
work on \emph{general mechanisms} for either finding differences in versions
or specifying differential comparison problems is relatively limited.

\subsection{Differential Testing (and Fuzzing) for Specialized Settings:}
McKeeman's paper on using differential testing to find bugs in
compilers is the earliest and most influential work in the
field~\cite{Differential}, and essentially formally defined the
general concept of differential testing, albeit in a single, specific domain.  Compiler testing (as well as other
language tool testing, e.g., for debuggers~\cite{lehmann2018feedback}
or static analysis tools~\cite{StaticDiff}) has a long and
distinguished line of work based in differential
approaches~\cite{barany2018finding,chen2019deep,chen2016coverage};
indeed, by and large differential testing is the only means for
finding wrong-compile bugs that is in widespread use, and is essential
to the success of tools such as Csmith~\cite{csmith}.  Klinger et al.~\cite{klinger2019differentially} similarly apply differential testing to check soundness and precision across program analyzers, an approach conceptually adjacent to our own differential evaluation of static analysis tools~\cite{StaticDiff}.

File systems
have also been a target for differential approaches~\cite{ICSEDiff}.
Recently machine learning systems have also been a focus of
differential
approaches~\cite{herbold2023differential,li2026enhancing}.  Security
and privacy testing also makes significant use of differential
approaches~\cite{jung2008privacy,chen2015guided}; a notable recent
example is the work of Nie et al.~\cite{nie2023coverage} that
discovered 11 novel root causes of X.509 certificate validation
divergences between SSL/TLS implementations, and demonstrated that RFC
5280 is fundamentally ambigious.  Differential approaches have also
been used to detect timing and other side-channel
leaks~\cite{nilizadeh2019diffuzz} and to expose platform-dependent
behavioral divergence in heterogeneous
applications~\cite{zhang2021heterofuzz}. Fine-grained differential
fuzzing has recently been pushed into JavaScript engines as well,
targeting semantic divergence below the level of crash
detection~\cite{wachter2025dumpling}.  All of these domains may
provide interesting case studies and benchmarks for this proposal.

Rigger, Ba, Su and others~\cite{Rigger2020PQS,Rigger2020TLP,Ba2024DQP} have
produced a substantial body of work on testing databases using
differential \emph{inputs} that could map to a differential fuzzing
problem, by extending our approach's insights to apply to cases where
part of a test problem is held constant (e.g., the database under test
plus a differential query plan pair) and another part (database
contents or database system configurations options) is varied.  We
believe that additional, perhaps surprising, problems that can be
fruitfully treated as differential
behavior discovery can be identified once differential fuzzing is made
highly effective and differential specification is made less onerous.
Here the underlying approaches are (in the case of the differentail
query plans paper) explicitly identified as ``differential'' but not
in the way usually considered in differential fuzzing.

The recent work of Jeong et al. on differential fuzzing of Ethereum
consensus clients is interesting for exploring both generation and
specification aspects of the problem~\cite{Jeong2026SpecTrum}.
Similarly, Yang et al.~\cite{yang2021consensus} used multi-transaction
differential fuzzing to find consensus bugs across Ethereum client
implementations, which form a particularly interesting setting where independently-implemented ``same functionality'' is security-critical.

\paragraph{Differential Testing, Fuzzing, and Specification as a First-class Problem:}
As noted, the most important such work, with respect to \emph{detection} of
differential bugs, is probably that of Petsios et
al. on the NEZHA framework~\cite{nezha}.  As discussed extensively
above, this work revitalized interest in the problem of working out
\emph{general purpose} methods of finding divergences between software
systems using state-of-the-art fuzzers.  While we hypothesize that
some algorithmic
fuzzing advances presented in the NEZHA work may be highly significant
only for
libFuzzer-like novelty functions, the general
framework for posing the problem and focusing on executions that can
reveal differences is a major precursor to this proposal's approach.

Böhme et al.~\cite{bohme2021residual} formalized residual risk and
discovery-probability estimation for greybox campaigns, and Baez et
al.~\cite{baez2025evt} extended this style of analysis specifically to
differential fuzzing using extreme value theory, giving a principled
way to estimate the probability that continued fuzzing will expose a
larger divergence.

Person et al.~\cite{Person2008DSE} explored using symbolic execution
to solve general differential testing problems, and while the proposed
methods have generally not been found to scale to most software
systems, the basic idea of characterizing behavioral differences
informs this proposal, including the idea of an AFL++ mode that
characterizes behavior changes as likely intended or unintended (based on both nature and
frequency of occurrence) based on large bodies of coverage-driven
exploration rather than symbolic execution.

PI Groce's previous work in
the TSTL testing framework for Python is, to our knowledge, the most
extensive previous effort to make differential specification an easier
problem, by mapping sequences of actions over two implementations into
each other, and allowing for use of abstraction functions over inputs
and outputs in such a setting~\cite{ISSTA15,tstl}.  However, the idea
of comparative or variant testing as critical to testing DSLs or
languages has been understood, if not heavily explored at the language
level, since early work on e.g., \emph{lava}~\cite{Sirer1999Lava}.

The Knight-Leveson work on
correlated failures in ``N-version'' software~\cite{KNIGHT1985159} famously suggested that
the fruits of differential testing might be less substantial than
anticipated, a problem for which we believe agentic/AI coding may now provide a
substantial solution.  Differential testing is one of the most
interesting solutions to the general problem of constructing
high-powered oracles for software testing~\cite{Barr2015}.
Metamorphic testing can be seen as an attempt to apply differential
testing without requiring actual multiple versions of
code~\cite{Segura2016}. The work of Rigger, Ba, and Su above on
databases is a form of metamorphic testing, for the most part.